\documentclass[10pt,a4paper,onecolumn]{article}
\usepackage{marginnote}
\usepackage{graphicx}
\usepackage{xcolor}
\usepackage{authblk,etoolbox}
\usepackage{titlesec}
\usepackage{calc}
\usepackage{tikz}
\usepackage{hyperref}
\hypersetup{colorlinks,breaklinks=true,
            urlcolor=[rgb]{0.0, 0.5, 1.0},
            linkcolor=[rgb]{0.0, 0.5, 1.0}}
\usepackage{caption}
\usepackage{tcolorbox}
\usepackage{amssymb,amsmath}
\usepackage{ifxetex,ifluatex}
\usepackage{seqsplit}
\usepackage{xstring}

\usepackage{float}
\let\origfigure\figure
\let\endorigfigure\endfigure
\renewenvironment{figure}[1][2] {
    \expandafter\origfigure\expandafter[H]
} {
    \endorigfigure
}


\usepackage{fixltx2e} % provides \textsubscript
\usepackage[
  backend=biber,
%  style=alphabetic,
%  citestyle=numeric
]{biblatex}
\bibliography{paper.bib}

% --- Splitting \texttt --------------------------------------------------

\let\textttOrig=\texttt
\def\texttt#1{\expandafter\textttOrig{\seqsplit{#1}}}
\renewcommand{\seqinsert}{\ifmmode
  \allowbreak
  \else\penalty6000\hspace{0pt plus 0.02em}\fi}


% --- Pandoc does not distinguish between links like [foo](bar) and
% --- [foo](foo) -- a simplistic Markdown model.  However, this is
% --- wrong:  in links like [foo](foo) the text is the url, and must
% --- be split correspondingly.
% --- Here we detect links \href{foo}{foo}, and also links starting
% --- with https://doi.org, and use path-like splitting (but not
% --- escaping!) with these links.
% --- Another vile thing pandoc does is the different escaping of
% --- foo and bar.  This may confound our detection.
% --- This problem we do not try to solve at present, with the exception
% --- of doi-like urls, which we detect correctly.


\makeatletter
\let\href@Orig=\href
\def\href@Urllike#1#2{\href@Orig{#1}{\begingroup
    \def\Url@String{#2}\Url@FormatString
    \endgroup}}
\def\href@Notdoi#1#2{\def\tempa{#1}\def\tempb{#2}%
  \ifx\tempa\tempb\relax\href@Urllike{#1}{#2}\else
  \href@Orig{#1}{#2}\fi}
\def\href#1#2{%
  \IfBeginWith{#1}{https://doi.org}%
  {\href@Urllike{#1}{#2}}{\href@Notdoi{#1}{#2}}}
\makeatother

\newlength{\cslhangindent}
\setlength{\cslhangindent}{1.5em}
\newlength{\csllabelwidth}
\setlength{\csllabelwidth}{3em}
\newenvironment{CSLReferences}[3] % #1 hanging-ident, #2 entry spacing
 {% don't indent paragraphs
  \setlength{\parindent}{0pt}
  % turn on hanging indent if param 1 is 1
  \ifodd #1 \everypar{\setlength{\hangindent}{\cslhangindent}}\ignorespaces\fi
  % set entry spacing
  \ifnum #2 > 0
  \setlength{\parskip}{#2\baselineskip}
  \fi
 }%
 {}
\usepackage{calc}
\newcommand{\CSLBlock}[1]{#1\hfill\break}
\newcommand{\CSLLeftMargin}[1]{\parbox[t]{\csllabelwidth}{#1}}
\newcommand{\CSLRightInline}[1]{\parbox[t]{\linewidth - \csllabelwidth}{#1}}
\newcommand{\CSLIndent}[1]{\hspace{\cslhangindent}#1}

% --- Page layout -------------------------------------------------------------
\usepackage[top=3.5cm, bottom=3cm, right=1.5cm, left=1.0cm,
            headheight=2.2cm, reversemp, includemp, marginparwidth=4.5cm]{geometry}

% --- Default font ------------------------------------------------------------
\renewcommand\familydefault{\sfdefault}

% --- Style -------------------------------------------------------------------
\renewcommand{\bibfont}{\small \sffamily}
\renewcommand{\captionfont}{\small\sffamily}
\renewcommand{\captionlabelfont}{\bfseries}

% --- Section/SubSection/SubSubSection ----------------------------------------
\titleformat{\section}
  {\normalfont\sffamily\Large\bfseries}
  {}{0pt}{}
\titleformat{\subsection}
  {\normalfont\sffamily\large\bfseries}
  {}{0pt}{}
\titleformat{\subsubsection}
  {\normalfont\sffamily\bfseries}
  {}{0pt}{}
\titleformat*{\paragraph}
  {\sffamily\normalsize}


% --- Header / Footer ---------------------------------------------------------
\usepackage{fancyhdr}
\pagestyle{fancy}
\fancyhf{}
%\renewcommand{\headrulewidth}{0.50pt}
\renewcommand{\headrulewidth}{0pt}
\fancyhead[L]{\hspace{-0.75cm}\includegraphics[width=5.5cm]{logo.png}}
\fancyhead[C]{}
\fancyhead[R]{}
\renewcommand{\footrulewidth}{0.25pt}

\fancyfoot[L]{\parbox[t]{0.98\headwidth}{\footnotesize{\sffamily Hemantakumar Phurailatpam, Otto Akseli Hannuksela, (2024). \texttt{gwsnr}:
A python package for efficient signal-to-noise calculation of
gravitational-waves. \textit{PendingJournal of Open Source Software}, Pending(Pending), Pending. \url{https://doi.org/10.xxxxxx/draft}}}}


\fancyfoot[R]{\sffamily \thepage}
\makeatletter
\let\ps@plain\ps@fancy
\fancyheadoffset[L]{4.5cm}
\fancyfootoffset[L]{4.5cm}

% --- Macros ---------

\definecolor{linky}{rgb}{0.0, 0.5, 1.0}

\newtcolorbox{repobox}
   {colback=red, colframe=red!75!black,
     boxrule=0.5pt, arc=2pt, left=6pt, right=6pt, top=3pt, bottom=3pt}

\newcommand{\ExternalLink}{%
   \tikz[x=1.2ex, y=1.2ex, baseline=-0.05ex]{%
       \begin{scope}[x=1ex, y=1ex]
           \clip (-0.1,-0.1)
               --++ (-0, 1.2)
               --++ (0.6, 0)
               --++ (0, -0.6)
               --++ (0.6, 0)
               --++ (0, -1);
           \path[draw,
               line width = 0.5,
               rounded corners=0.5]
               (0,0) rectangle (1,1);
       \end{scope}
       \path[draw, line width = 0.5] (0.5, 0.5)
           -- (1, 1);
       \path[draw, line width = 0.5] (0.6, 1)
           -- (1, 1) -- (1, 0.6);
       }
   }

% --- Title / Authors ---------------------------------------------------------
% patch \maketitle so that it doesn't center
\patchcmd{\@maketitle}{center}{flushleft}{}{}
\patchcmd{\@maketitle}{center}{flushleft}{}{}
% patch \maketitle so that the font size for the title is normal
\patchcmd{\@maketitle}{\LARGE}{\LARGE\sffamily}{}{}
% patch the patch by authblk so that the author block is flush left
\def\maketitle{{%
  \renewenvironment{tabular}[2][]
    {\begin{flushleft}}
    {\end{flushleft}}
  \AB@maketitle}}
\makeatletter
\renewcommand\AB@affilsepx{ \protect\Affilfont}
%\renewcommand\AB@affilnote[1]{{\bfseries #1}\hspace{2pt}}
\renewcommand\AB@affilnote[1]{{\bfseries #1}\hspace{3pt}}
\renewcommand{\affil}[2][]%
   {\newaffiltrue\let\AB@blk@and\AB@pand
      \if\relax#1\relax\def\AB@note{\AB@thenote}\else\def\AB@note{#1}%
        \setcounter{Maxaffil}{0}\fi
        \begingroup
        \let\href=\href@Orig
        \let\texttt=\textttOrig
        \let\protect\@unexpandable@protect
        \def\thanks{\protect\thanks}\def\footnote{\protect\footnote}%
        \@temptokena=\expandafter{\AB@authors}%
        {\def\\{\protect\\\protect\Affilfont}\xdef\AB@temp{#2}}%
         \xdef\AB@authors{\the\@temptokena\AB@las\AB@au@str
         \protect\\[\affilsep]\protect\Affilfont\AB@temp}%
         \gdef\AB@las{}\gdef\AB@au@str{}%
        {\def\\{, \ignorespaces}\xdef\AB@temp{#2}}%
        \@temptokena=\expandafter{\AB@affillist}%
        \xdef\AB@affillist{\the\@temptokena \AB@affilsep
          \AB@affilnote{\AB@note}\protect\Affilfont\AB@temp}%
      \endgroup
       \let\AB@affilsep\AB@affilsepx
}
\makeatother
\renewcommand\Authfont{\sffamily\bfseries}
\renewcommand\Affilfont{\sffamily\small\mdseries}
\setlength{\affilsep}{1em}


\ifnum 0\ifxetex 1\fi\ifluatex 1\fi=0 % if pdftex
  \usepackage[T1]{fontenc}
  \usepackage[utf8]{inputenc}

\else % if luatex or xelatex
  \ifxetex
    \usepackage{mathspec}
    \usepackage{fontspec}

  \else
    \usepackage{fontspec}
  \fi
  \defaultfontfeatures{Ligatures=TeX,Scale=MatchLowercase}

\fi
% use upquote if available, for straight quotes in verbatim environments
\IfFileExists{upquote.sty}{\usepackage{upquote}}{}
% use microtype if available
\IfFileExists{microtype.sty}{%
\usepackage{microtype}
\UseMicrotypeSet[protrusion]{basicmath} % disable protrusion for tt fonts
}{}

\usepackage{hyperref}
\hypersetup{unicode=true,
            pdftitle={gwsnr: A python package for efficient signal-to-noise calculation of gravitational-waves},
            pdfborder={0 0 0},
            breaklinks=true}
\urlstyle{same}  % don't use monospace font for urls

% --- We redefined \texttt, but in sections and captions we want the
% --- old definition
\let\addcontentslineOrig=\addcontentsline
\def\addcontentsline#1#2#3{\bgroup
  \let\texttt=\textttOrig\addcontentslineOrig{#1}{#2}{#3}\egroup}
\let\markbothOrig\markboth
\def\markboth#1#2{\bgroup
  \let\texttt=\textttOrig\markbothOrig{#1}{#2}\egroup}
\let\markrightOrig\markright
\def\markright#1{\bgroup
  \let\texttt=\textttOrig\markrightOrig{#1}\egroup}


\usepackage{graphicx,grffile}
\makeatletter
\def\maxwidth{\ifdim\Gin@nat@width>\linewidth\linewidth\else\Gin@nat@width\fi}
\def\maxheight{\ifdim\Gin@nat@height>\textheight\textheight\else\Gin@nat@height\fi}
\makeatother
% Scale images if necessary, so that they will not overflow the page
% margins by default, and it is still possible to overwrite the defaults
% using explicit options in \includegraphics[width, height, ...]{}
\setkeys{Gin}{width=\maxwidth,height=\maxheight,keepaspectratio}
\IfFileExists{parskip.sty}{%
\usepackage{parskip}
}{% else
\setlength{\parindent}{0pt}
\setlength{\parskip}{6pt plus 2pt minus 1pt}
}
\setlength{\emergencystretch}{3em}  % prevent overfull lines
\providecommand{\tightlist}{%
  \setlength{\itemsep}{0pt}\setlength{\parskip}{0pt}}
\setcounter{secnumdepth}{0}
% Redefines (sub)paragraphs to behave more like sections
\ifx\paragraph\undefined\else
\let\oldparagraph\paragraph
\renewcommand{\paragraph}[1]{\oldparagraph{#1}\mbox{}}
\fi
\ifx\subparagraph\undefined\else
\let\oldsubparagraph\subparagraph
\renewcommand{\subparagraph}[1]{\oldsubparagraph{#1}\mbox{}}
\fi

\title{\texttt{gwsnr}: A python package for efficient signal-to-noise
calculation of gravitational-waves}

        \author[1]{Hemantakumar Phurailatpam}
          \author[1]{Otto Akseli HANNUKSELA}
    
      \affil[1]{Department of Physics, The Chinese University of Hong
Kong, Shatin, New Territories, Hong Kong}
  \date{\vspace{-7ex}}

\begin{document}
\maketitle

\marginpar{

  \begin{flushleft}
  %\hrule
  \sffamily\small

  {\bfseries DOI:} \href{https://doi.org/10.xxxxxx/draft}{\color{linky}{10.xxxxxx/draft}}

  \vspace{2mm}

  {\bfseries Software}
  \begin{itemize}
    \setlength\itemsep{0em}
    \item \href{}{\color{linky}{Review}} \ExternalLink
    \item \href{https://github.com/hemantaph/ler}{\color{linky}{Repository}} \ExternalLink
    \item \href{https://zenodo.org/badge/latestdoi/626733473}{\color{linky}{Archive}} \ExternalLink
  \end{itemize}

  \vspace{2mm}

  \par\noindent\hrulefill\par

  \vspace{2mm}

  {\bfseries Editor:} \href{}{Pending Editor} \ExternalLink \\
  \vspace{1mm}
    {\bfseries Reviewers:}
  \begin{itemize}
  \setlength\itemsep{0em}
    \item \href{https://github.com/Pending Reviewer}{@Pending Reviewer}
    \item \href{https://github.com/}{@}
    \end{itemize}
    \vspace{2mm}

  {\bfseries Submitted:} 13th Dec 2024\\
  {\bfseries Published:} 

  \vspace{2mm}
  {\bfseries License}\\
  Authors of papers retain copyright and release the work under a Creative Commons Attribution 4.0 International License (\href{http://creativecommons.org/licenses/by/4.0/}{\color{linky}{CC BY 4.0}}).

  
  \end{flushleft}
}

\section{Summary}\label{summary}

Gravitational waves (GWs), ripples in spacetime predicted by Einstein's
theory of General Relativity, have revolutionized astrophysics since
their first detection in 2015 (Abbott, B. P. et al. (2016a), B.  P.
Abbott et al. (2016)). These waves are emitted by cataclysmic events
like the merging of binary black holes (BBHs), binary neutron stars
(BNSs) and BH-NS pairs, providing a unique window into the cosmos. A
critical aspect of GW analysis is the Signal-to-Noise Ratio (SNR). SNR
quantifies the strength of a GW signal relative to the background noise
in a detector, like LIGO (The LIGO Scientific Collaboration et al.
(2015), B. P. Abbott et al. (2020), Buikema et al. (2020)), Virgo (F.
Acernese et al. (2014), F. Acernese et al. (2019)) or KAGRA (Akutsu et
al. (2020), Aso et al. (2013)). This ratio is pivotal in confirming the
detection of GWs and extracting astrophysical information from them
(Abbott, B. P. et al. (2016b)). However, specific scenarios in GW
research, particularly in simulations of detectable GW events (B. P.
Abbott et al. (2016)) and in hierarchical Bayesian analysis (Thrane and
Talbot (2019)) where selection effects are considered, demand extensive
and efficient computation of SNR. This requirement presents a
significant challenge, as conventional computational approaches, such as
noise-weighted inner product, are typically time-consuming and
impractical for such specialized and large-scale analyses (Taylor and
Gerosa (2018), Gerosa et al. (2020)).

\section{Statement of Need}\label{statement-of-need}

The \emph{\texttt{gwsnr}} Python package addresses the need for
efficient SNR computation in GW research. It provides a flexible and
user-friendly interface, allowing users to combine various detector
noise models, waveform models, detector configurations, and signal
parameters. \emph{\texttt{gwsnr}} enhances SNR calculations through
several key features. Firstly, it utilizes an innovative interpolation
method, employing a partial-scaling approach for accurately
interpolating the SNR of GWs from spin-less and spin-aligned binary
systems. Secondly, the package features a noise-weighted inner product
method, similar to that in the \emph{\texttt{bilby}} package (Ashton et
al. (2019), Ashton, Gregory et al. (2022)), but enhanced with
multiprocessing capabilities. This parallel processing is crucial for
handling large datasets and computationally intensive analyses. Thirdly,
a trained Artificial Neural Network (ANN) model is incorporated for
rapid `probability of detection' (Pdet) estimation for BBH systems with
spin precession. Lastly, \emph{\texttt{gwsnr}} leverages
\emph{\texttt{numpy}} (NumPy Community (2022)) vectorization, and
\emph{\texttt{numba}}'s (Lam, Pitrou, and Seibert (2022)) and
\emph{\texttt{JAX}}'s (James Bradbury and others (2018)) Just-In-Time
compiler (\texttt{numbba.njit} and \texttt{jax.jit}), which optimizes
performance by compiling Python code into machine code at runtime,
drastically reducing execution times. This combination of advanced
techniques and user-friendly design makes gwsnr a valuable tool for GW
data analysis, particularly in simulating detectable compact binary
mergers, determining rates of both lensed and unlensed GW events (as
demonstrated by its use in the \emph{\texttt{ler}} package; Phurailatpam
et al. (2024), Ng et al. (2024), More and Phurailatpam (2025), Janquart
et al. (2023), R. Abbott et al. (2021), Collaboration et al. (2023),
Wierda et al. (2021), Wempe et al. (2022)), and will help in the
analysis of selection effects within hierarchical Bayesian frameworks
(Thrane and Talbot (2019)).

\section{Mathematical Formulation}\label{mathematical-formulation}

The \texttt{gwsnr} package provides two efficient methods for computing
the optimal SNR in GW data analysis: the Noise-Weighted Inner Product
Method with Multiprocessing and the Partial Scaling Interpolation
Method. In addition, there are two approaches for estimating
\(P_{\rm det}\) for precessing systems: ANN-based \(P_{\rm det}\)
Estimation and the Partial Scaling Interpolation Method with SNR
recalculation. Extensive details of these methods can be found in the
package documentation (Phurailatpam and Hannuksela (2025)).

\subsubsection{Noise-Weighted Inner Product Method with
Multiprocessing}\label{noise-weighted-inner-product-method-with-multiprocessing}

The noise-weighted inner product is a robust and widely used technique,
suitable for any frequency-domain gravitational waveform, including
complex models with spin precession and higher-order harmonics available
in \texttt{lalsimulation} (LIGO Scientific Collaboration, Virgo
Collaboration, and KAGRA Collaboration (2018)). Following (Allen et al.
(2012)), the inner product between two frequency-domain signals,
\(\tilde{a}(f)\) and \(\tilde{b}(f)\), is defined as:

\[
\left< a | b \right> = 4 \Re \int_{f_{\min}}^{f_{\max}} \frac{\tilde{a}(f)\tilde{b}^*(f)}{S_n(f)} df
\]

Here, \(S_n(f)\) is the one-sided power spectral density of the detector
noise, and \((f_{\min}, f_{\max})\) is the analysis frequency band. The
optimal SNR \(\rho\), is the norm of the inner-product for the given
signal \(h\) : \(\rho = \sqrt{\langle h | h \rangle}\). For a
gravitational wave signal composed of plus (\(h_+\)) and cross
(\(h_\times\)) polarizations, and assuming orthogonality between them,
the SNR can be expressed in terms of the detector's antenna patterns,
\(F_+\) and \(F_\times\):

\[
\rho = \sqrt{ F_+^2 \left< \tilde{h}_+ | \tilde{h}_+ \right> + F_{\times}^2 \left< \tilde{h}_{\times} | \tilde{h}_{\times} \right> }
\]

While this approach is versatile, it can be computationally intensive,
with waveform generation representing the primary bottleneck. The
\texttt{gwsnr} package addresses this challenge by parallelizing
waveform generation across multiple CPU cores and accelerating the
antenna pattern and inner product calculations using \texttt{numba.njit}
compilation. Additionally, \texttt{gwsnr} provides optional support for
JAX-based waveform generation and acceleration via the \texttt{ripple}
waveform library (Edwards et al. 2024), utilizing \texttt{jax.jit} for
just-in-time compilation and \texttt{jax.vmap} for efficient batched
operations.

\subsubsection{Partial Scaling Interpolation
Method}\label{partial-scaling-interpolation-method}

For non-spinning or aligned-spin binary systems restricted to the
dominant harmonic mode, \texttt{gwsnr} implements a highly efficient
interpolation-based technique called the Partial Scaling method. This
approach, adapted from the FINDCHIRP algorithm (Allen et al. (2012)),
decouples the computationally expensive parts of the SNR calculation
from the extrinsic source parameters. It achieves this by defining a
``partial-scaled SNR'' \(\rho_{1/2}\), which isolates the dependence on
the intrinsic parameters (masses and spins). For a given full IMR
waveform SNR, \(\rho_{\text{full}}\), the partial SNR is defined as:

\[
\rho_{1/2} = \left(\frac{D_\mathrm{eff}}{1~\mathrm{Mpc}}\right) \left(\frac{\mathcal{M}}{M_\odot}\right)^{-5/6} \times \rho_{\text{full}}
\]

Here, \(\mathcal{M}\) is the chirp mass and \(D_{\text{eff}}\) is the
effective distance, which encapsulates the luminosity distance, sky
location, and detector orientation wrt the binary. Since \(\rho_{1/2}\)
depends only on the intrinsic properties of the binary, its value can be
pre-computed on a grid and stored. For non-spinning systems, this is a
two-dimensional grid of total mass (\(M\)) and mass ratio (\(q\)), while
for aligned-spin systems, it is a four-dimensional grid that also
includes the two spin magnitudes. To find the SNR for a new binary,
\texttt{gwsnr} performs a rapid cubic spline interpolation on the
pre-computed grid to find the corresponding \(\rho_{1/2}\) value. The
final SNR is then recovered almost instantaneously by applying the
scaling transformation:

\[
\rho = \rho_{1/2} \times \left(\frac{\mathcal{M}}{M_\odot}\right)^{5/6} \times \left(\frac{1~\mathrm{Mpc}}{D_\mathrm{eff}}\right)
\]

This procedure transforms a computationally intensive integration into a
simple, JIT-compiled interpolation and multiplication, enabling massive
performance gains for large-scale population studies.

\subsubsection{ANN-based Pdet
Estimation}\label{ann-based-pdet-estimation}

The \texttt{gwsnr} package now incorporates an artificial neural network
(ANN) model, developed using TensorFlow (Abadi et al. (2015)) and
scikit-learn (Pedregosa et al. (2011)), to rapidly estimate
\(P_{\rm det}\) in binary black hole (BBH) systems using the
IMRPhenomXPHM waveform approximant. This complex IMR waveform model
accounts for spin-precessing systems with subdominant harmonics. The ANN
model is especially useful when precise signal-to-noise ratio (SNR)
calculations are not critical, providing a quick and effective means of
estimating \(P_{\rm det}\). This value indicates detectability under
Gaussian noise by determining if the SNR exceeds a certain threshold
(e.g., \(\rho_{\rm th}=8\)). Trained on a large dataset from the
\texttt{ler} package, the ANN model uses `partial scaled SNR' values as
a primary input, reducing input dimensionality from 15 to 5 and
enhancing accuracy. This approach offers a practical solution for
assessing detectability under specified conditions. Other similar
efforts with ANN models are detailed in (Chapman-Bird et al. (2023),
Gerosa et al. (2020), Callister et al. (2024)).

In addition to providing trained ANN models for specific configurations,
\texttt{gwsnr} offers users the flexibility to develop and train custom
models tailored to their unique requirements. This adaptability allows
for optimization based on variations in detector sensitivity,
gravitational-wave properties, and other research-specific factors,
ensuring maximum model effectiveness across different scenarios.

\subsubsection{Partial Scaling Interpolation Method with SNR
Recalculation for Pdet
Estimation}\label{partial-scaling-interpolation-method-with-snr-recalculation-for-pdet-estimation}

While the Partial Scaling method is highly efficient for aligned-spin
systems, its utility can be further enhanced by recalculating the SNR
for precessing systems within a predefined small range of generated
SNRs. This is done by first obtaining optimal SNRs with the Partial
Scaling method, selecting the SNRs near \(\rho_{\rm th}\), and then
recalculating the SNRs for these systems using the Noise-Weighted Inner
Product Method. This approach allows us to leverage the speed of the
Partial Scaling method while ensuring accurate SNR values for systems
close to the detection threshold. The recalculated SNRs can then be used
to estimate \(P_{\rm det}\), providing a balance between computational
efficiency and accuracy.

\section{Acknowledgements}\label{acknowledgements}

The author gratefully acknowledges the substantial contributions from
all who supported this research. Special thanks go to my academic
advisors for their invaluable guidance and unwavering support. The
interactions with my research colleagues have greatly enriched this
work. The Department of Physics at The Chinese University of Hong Kong
is acknowledged for the Postgraduate Studentship that made this research
possible. Thanks are also due to the LIGO Laboratory for the
computational resources, supported by National Science Foundation Grants
No.~PHY-0757058 and No.~PHY-0823459.

\section{References}\label{references}

\phantomsection\label{refs}
\begin{CSLReferences}{1}{0}
\bibitem[\citeproctext]{ref-tensorflow:2015}
Abadi, Martín, Ashish Agarwal, Paul Barham, Eugene Brevdo, Zhifeng Chen,
Craig Citro, Greg S. Corrado, et al. 2015. {``{TensorFlow}: Large-Scale
Machine Learning on Heterogeneous Systems.''}
\url{https://www.tensorflow.org/}.

\bibitem[\citeproctext]{ref-Abbott:2016:rates}
Abbott, B. P., R. Abbott, T. D. Abbott, M. R. Abernathy, F. Acernese, K.
Ackley, C. Adams, et al. 2016. {``ASTROPHYSICAL IMPLICATIONS OF THE
BINARY BLACK HOLE MERGER GW150914.''} \emph{The Astrophysical Journal
Letters} 818 (2): L22.
\url{https://doi.org/10.3847/2041-8205/818/2/l22}.

\bibitem[\citeproctext]{ref-Abbott:2016}
Abbott, B. P., R. Abbott, T. D. Abbott, M. R. Abernathy, F. Acernese, K.
Ackley, C. Adams, et al. 2016a. {``Observation of Gravitational Waves
from a Binary Black Hole Merger.''} \emph{Physical Review Letters} 116
(6). \url{https://doi.org/10.1103/physrevlett.116.061102}.

\bibitem[\citeproctext]{ref-Abbott:2016:detection}
Abbott, B. P., R. Abbott, T. D. Abbott, M. R. Abernathy, F. Acernese, K.
Ackley, C. Adams, et al. 2016b. {``GW150914: First Results from the
Search for Binary Black Hole Coalescence with Advanced LIGO.''}
\emph{Physical Review D} 93 (12).
\url{https://doi.org/10.1103/physrevd.93.122003}.

\bibitem[\citeproctext]{ref-Abbott:2016:pe}
Abbott, B. P., R. Abbott, T. D. Abbott, M. R. Abernathy, F. Acernese, K.
Ackley, C. Adams, et al. 2016. {``Properties of the Binary Black Hole
Merger GW150914.''} \emph{Physical Review Letters} 116 (24).
\url{https://doi.org/10.1103/physrevlett.116.241102}.

\bibitem[\citeproctext]{ref-Abbott:2020}
Abbott, B. P., R. Abbott, T. D. Abbott, S. Abraham, F. Acernese, K.
Ackley, C. Adams, et al. 2020. {``Prospects for Observing and Localizing
Gravitational-Wave Transients with Advanced LIGO, Advanced Virgo and
KAGRA.''} \emph{Living Reviews in Relativity} 23 (1).
\url{https://doi.org/10.1007/s41114-020-00026-9}.

\bibitem[\citeproctext]{ref-Abbott:2021}
Abbott, R., T. D. Abbott, S. Abraham, F. Acernese, K. Ackley, A. Adams,
C. Adams, et al. 2021. {``Search for Lensing Signatures in the
Gravitational-Wave Observations from the First Half of LIGO--Virgo's
Third Observing Run.''} \emph{The Astrophysical Journal} 923 (1): 14.
\url{https://doi.org/10.3847/1538-4357/ac23db}.

\bibitem[\citeproctext]{ref-VIRGO:2019}
Acernese, F., M. Agathos, L. Aiello, A. Allocca, A. Amato, S. Ansoldi,
S. Antier, et al. 2019. {``Increasing the Astrophysical Reach of the
Advanced Virgo Detector via the Application of Squeezed Vacuum States of
Light.''} \emph{Phys. Rev. Lett.} 123 (December): 231108.
\url{https://doi.org/10.1103/PhysRevLett.123.231108}.

\bibitem[\citeproctext]{ref-VIRGO:2015}
Acernese, F, M Agathos, K Agatsuma, D Aisa, N Allemandou, A Allocca, J
Amarni, et al. 2014. {``Advanced Virgo: A Second-Generation
Interferometric Gravitational Wave Detector.''} \emph{Classical and
Quantum Gravity} 32 (2): 024001.
\url{https://doi.org/10.1088/0264-9381/32/2/024001}.

\bibitem[\citeproctext]{ref-Akutsu:2020}
Akutsu, T., M. Ando, K. Arai, Y. Arai, S. Araki, A. Araya, N. Aritomi,
et al. 2020. {``Overview of KAGRA: Detector Design and Construction
History.''} \url{https://arxiv.org/abs/2005.05574}.

\bibitem[\citeproctext]{ref-Allen:2012}
Allen, Bruce, Warren G. Anderson, Patrick R. Brady, Duncan A. Brown, and
Jolien D. E. Creighton. 2012. {``FINDCHIRP: An Algorithm for Detection
of Gravitational Waves from Inspiraling Compact Binaries.''}
\emph{Physical Review D} 85 (12).
\url{https://doi.org/10.1103/physrevd.85.122006}.

\bibitem[\citeproctext]{ref-Ashton:2022}
Ashton, Gregory, Sylvia Biscoveanu, Neil Cornish, Isaac Dal Canto,
Prayush Kumar, Duncan Meacher, Hannah Middleton, Divyansh Mistry, Rory
Smith, and Tom Stevenson. 2022. {``{bilby: a user-friendly Bayesian
inference library}.''} \emph{GitHub Repository}. GitHub.
\url{https://github.com/GregoryAshton/Bilby}.

\bibitem[\citeproctext]{ref-Ashton:2019}
Ashton, Gregory, Moritz Hübner, Paul D. Lasky, Colm Talbot, Kendall
Ackley, Sylvia Biscoveanu, Qi Chu, et al. 2019. {``Bilby: A
User-Friendly Bayesian Inference Library for Gravitational-Wave
Astronomy.''} \emph{The Astrophysical Journal Supplement Series} 241
(2): 27. \url{https://doi.org/10.3847/1538-4365/ab06fc}.

\bibitem[\citeproctext]{ref-Aso:2013}
Aso, Yoichi, Yuta Michimura, Kentaro Somiya, Masaki Ando, Osamu
Miyakawa, Takanori Sekiguchi, Daisuke Tatsumi, and Hiroaki Yamamoto.
2013. {``Interferometer Design of the KAGRA Gravitational Wave
Detector.''} \emph{Phys. Rev. D} 88 (August): 043007.
\url{https://doi.org/10.1103/PhysRevD.88.043007}.

\bibitem[\citeproctext]{ref-Buikema:2020}
Buikema, A., C. Cahillane, G. L. Mansell, C. D. Blair, R. Abbott, C.
Adams, R. X. Adhikari, et al. 2020. {``Sensitivity and Performance of
the Advanced LIGO Detectors in the Third Observing Run.''} \emph{Phys.
Rev. D} 102 (September): 062003.
\url{https://doi.org/10.1103/PhysRevD.102.062003}.

\bibitem[\citeproctext]{ref-Callister:2024}
Callister, Thomas A. et al. 2024. {``Neural Network Emulator of the
Advanced LIGO and Advanced Virgo Selection Function.''} \emph{Physical
Review D} 110 (12). \url{https://doi.org/10.1103/physrevd.110.123041}.

\bibitem[\citeproctext]{ref-ChapmanBird:2023}
Chapman-Bird, Christian E A et al. 2023. {``Rapid Determination of LISA
Sensitivity to Extreme Mass Ratio Inspirals with Machine Learning.''}
\emph{Monthly Notices of the Royal Astronomical Society} 522 (4):
6043--54. \url{https://doi.org/10.1093/mnras/stad1397}.

\bibitem[\citeproctext]{ref-ligolensing:2023}
Collaboration, The LIGO Scientific, the Virgo Collaboration, the KAGRA
Collaboration, R. Abbott, H. Abe, F. Acernese, K. Ackley, et al. 2023.
{``Search for Gravitational-Lensing Signatures in the Full Third
Observing Run of the LIGO-Virgo Network.''}
\url{https://arxiv.org/abs/2304.08393}.

\bibitem[\citeproctext]{ref-Edwards:2023}
Edwards, Thomas D. P., Kaze W. K. Wong, Kelvin K. H. Lam, Adam Coogan,
Daniel Foreman-Mackey, Maximiliano Isi, and Aaron Zimmerman. 2024.
{``{Differentiable and hardware-accelerated waveforms for gravitational
wave data analysis}.''} \emph{Phys. Rev. D} 110 (6): 064028.
\url{https://doi.org/10.1103/PhysRevD.110.064028}.

\bibitem[\citeproctext]{ref-Gerosa:2020}
Gerosa, Davide et al. 2020. {``Gravitational-Wave Selection Effects
Using Neural-Network Classifiers.''} \emph{Physical Review D} 102 (10).
\url{https://doi.org/10.1103/physrevd.102.103020}.

\bibitem[\citeproctext]{ref-jax:2018}
James Bradbury, Peter Hawkins, Roy Frostig, and Various others. 2018.
{``JAX: Composable Transformations of Python+NumPy Programs.''} GitHub.
\url{https://github.com/google/jax}.

\bibitem[\citeproctext]{ref-Janquart:2023}
Janquart, J, M Wright, S Goyal, J C L Chan, A Ganguly, Á Garrón, D
Keitel, et al. 2023. {``{Follow-up analyses to the O3 LIGO--Virgo--KAGRA
lensing searches}.''} \emph{Monthly Notices of the Royal Astronomical
Society} 526 (3): 3832--60.
\url{https://doi.org/10.1093/mnras/stad2909}.

\bibitem[\citeproctext]{ref-numba:2022}
Lam, Stan, Stéphane Pitrou, and Mark Seibert. 2022. {``Numba: A High
Performance Python Compiler.''} \emph{Numba Documentation}. Anaconda,
Inc. \url{https://numba.pydata.org/}.

\bibitem[\citeproctext]{ref-lalsuite:2018}
LIGO Scientific Collaboration, Virgo Collaboration, and KAGRA
Collaboration. 2018. {``{LVK} {A}lgorithm {L}ibrary - {LALS}uite.''}
Free software (GPL). \url{https://doi.org/10.7935/GT1W-FZ16}.

\bibitem[\citeproctext]{ref-More:2025}
More, Anupreeta, and Hemantakumar Phurailatpam. 2025. {``Gravitational
Lensing: Towards Combining the Multi-Messengers.''}
\url{https://arxiv.org/abs/2502.02536}.

\bibitem[\citeproctext]{ref-Leo:2024}
Ng, Leo C. Y., Justin Janquart, Hemantakumar Phurailatpam, Harsh Narola,
Jason S. C. Poon, Chris Van Den Broeck, and Otto A. Hannuksela. 2024.
{``Uncovering Faint Lensed Gravitational-Wave Signals and Reprioritizing
Their Follow-up Analysis Using Galaxy Lensing Forecasts with Detected
Counterparts.''} \url{https://arxiv.org/abs/2403.16532}.

\bibitem[\citeproctext]{ref-numpy:2022}
NumPy Community. 2022. {``NumPy: A Fundamental Package for Scientific
Computing with Python.''} \emph{NumPy Website}. NumPy.
\url{https://numpy.org/}.

\bibitem[\citeproctext]{ref-scikitlearn:2011}
Pedregosa, Fabian, Gaël Varoquaux, Alexandre Gramfort, Vincent Michel,
Bertrand Thirion, Olivier Grisel, Mathieu Blondel, et al. 2011.
{``Scikit-Learn: Machine Learning in {P}ython.''} \emph{Journal of
Machine Learning Research} 12: 2825--30.

\bibitem[\citeproctext]{ref-gwsnr:documentation}
Phurailatpam, Hemantakumar, and Otto Akseli Hannuksela. 2025. {``Gwsnr:
Gravitational Wave Signal-to-Noise Ratio Computation Package
Documentation.''} \url{https://gwsnr.readthedocs.io/en/latest/}.

\bibitem[\citeproctext]{ref-ler:2024}
Phurailatpam, Hemantakumar, Anupreeta More, Harsh Narola, Ng Chung Yin,
Justin Janquart, Chris Van Den Broeck, Otto Akseli Hannuksela, Neha
Singh, and David Keitel. 2024. {``Ler : LVK (LIGO-Virgo-KAGRA
Collaboration) Event (Compact-Binary Mergers) Rate Calculator and
Simulator.''} \url{https://arxiv.org/abs/2407.07526}.

\bibitem[\citeproctext]{ref-Taylor:2018}
Taylor, Stephen R., and Davide Gerosa. 2018. {``Mining
Gravitational-Wave Catalogs to Understand Binary Stellar Evolution: A
New Hierarchical Bayesian Framework.''} \emph{Physical Review D} 98 (8).
\url{https://doi.org/10.1103/physrevd.98.083017}.

\bibitem[\citeproctext]{ref-LIGO:2015}
The LIGO Scientific Collaboration, J Aasi, B P Abbott, R Abbott, T
Abbott, M R Abernathy, K Ackley, et al. 2015. {``Advanced LIGO.''}
\emph{Classical and Quantum Gravity} 32 (7): 074001.
\url{https://doi.org/10.1088/0264-9381/32/7/074001}.

\bibitem[\citeproctext]{ref-Thrane:2019}
Thrane, Eric, and Colm Talbot. 2019. {``An Introduction to Bayesian
Inference in Gravitational-Wave Astronomy: Parameter Estimation, Model
Selection, and Hierarchical Models.''} \emph{Publications of the
Astronomical Society of Australia} 36.
\url{https://doi.org/10.1017/pasa.2019.2}.

\bibitem[\citeproctext]{ref-Wempe:2022}
Wempe, Ewoud, Léon V. E. Koopmans, A. Renske A. C. Wierda, Otto Akseli
Hannuksela, Alberto Agnello, Cyril Bonvin, Bendetta Bucciarelli, et al.
2022. {``A Lensing Multi-Messenger Channel: Combining LIGO-Virgo-Kagra
Lensed Gravitational-Wave Measurements with Euclid Observations.''}
\url{https://arxiv.org/abs/2204.08732}.

\bibitem[\citeproctext]{ref-Wierda:2021}
Wierda, A. Renske A. C., Ewoud Wempe, Otto A. Hannuksela, Léon V. E.
Koopmans, Alberto Agnello, Cyril Bonvin, Bendetta Bucciarelli, et al.
2021. {``Beyond the Detector Horizon: Forecasting Gravitational-Wave
Strong Lensing.''} \emph{The Astrophysical Journal} 921 (1): 154.
\url{https://doi.org/10.3847/1538-4357/ac1bb4}.

\end{CSLReferences}

\end{document}